\scanpage{54}
\mspage{50}
% archive sheet number: 53
$\mathrm{Fig}\,\mathrm{spat}\,\text{\unsure{éle|él|él.}}(\Sigma)$\inserted{ou $\mathrm{AtSpat}(\Sigma)$} l'ensemble des figures spatiales \uline{élémentaires} (i.e. telles que $\unsure{\Phi|\varPhi}$ ait un plus grand élément). On va mettre une structure d'atelier sur $\unsure{\mathrm{At}|\mathrm{AT}}(\Sigma)$, ainsi.
\[
\deleted{(\Phi,\rho)\trianglelefteq\Phi'}\qquad
\mathrm{AtSpat}(\Sigma)\subset\mathrm{Atens}(\Sigma)
\]
1) La relation $\trianglelefteq$ dans $\mathrm{Atspat}(\Sigma)$ est induite par la relation similaire dans $\mathrm{Atens}(\Sigma)$.

2) La relation $\mathrel{\overset{\unsure{\circ|o}}{<}\mkern-4mu<}$ dans $\mathrm{Atspat}\,\Sigma$ est \deleted{induite}\inserted{définie} comme plus haut : $(\Phi,\rho)\mathrel{\overset{\circ}{<}\mkern-4mu<}(\Phi',\rho')$ \unsure{sss|ssi} \deleted{par la relation similaire dans $\mathrm{Atens}\,\Sigma$} $\forall X\in\Phi,\ \exists X'=\psi(X)\in\Phi'$, avec $X\leq\deleted{\unsure{\psi}} \unsure{X'}\deleted{\unsure{(X)}}$ (NB. \deleted{(dans \unsure{Atens} pour $\unsure{\ll|\trianglelefteq}$, comme on voit} cet $X'$ est unique), et $\psi:\Phi\longrightarrow\Phi'$ est croissante \deleted{facilement \illegible{1} \unsure{constaté|constitué})} et transforme plus \unsure{grd|grand} \unsure{élémt|élément} en \unsure{idem|idm}.

3) La relation $\mathrel{|\unsure{o|\circ}|}$ dans $\mathrm{Atspat}(\Sigma)$ est \deleted{par contre} plus stricte que dans $\mathrm{Atens}(\Sigma)$ : on pose $(\Phi,\rho)\mathrel{|o|}(\Phi',\rho')$, sss (posant $X={}$plus grand \unsure{élémt|élément} de $\Phi$, et idem pour $\unsure{\alpha'|X'}$) non seulement on a $X\ne X'$ (relation $\mathrel{|o|}$ dans $\mathrm{Atens}$) mais même $X\mathrel{|o|}X'$ (ce qui est plus fort, vu que $X,\unsure{x'|X'}\ne0_\Sigma$).

Je dis que pour ces relations, on a \uline{un atelier}, \deleted{spatial} et même (je présume) un \uline{atelier} \uline{spatial}.
